% !TeX root = ../教材题库_全部条目_转写工作稿.tex
% Source: XJ-B1-C05-EX-P192-Y0318
% 第5章 / 5.4 函数 y=A sin(ωx+φ) 的图象与性质 / 教材 p.192 / PDF 第199页
% OCR未校对：公式与文字均可能有误；未完成数学TeX转写；禁止直接发给学生
% 原题及相邻原解截图（核对用）：../题目原图/XJ-B1-C05-EX-P192-Y0318-01.webp
% 本文件只在转写完成后作为题干使用；OCR 条目尚不是可靠题干。
\begin{bankitem}{5.4 例题 例题4}{第5章 / 5.4 函数 y=A sin(\ensuremath{\omega}x+\ensuremath{\varphi}) 的图象与性质 / 教材 p.192 / PDF 第199页}{XJ-B1-C05-EX-P192-Y0318}{OCR未校对：公式与文字均可能有误；未完成数学TeX转写；禁止直接发给学生}
\begin{ocrdraft}
[教材p.192 / PDF第199页]
\par
0二4]在同一平面直角坐标系中画出y=sin工，3-艹工十
\par
3吓=Sln工
\par
4
\par
在一个周期内的图象，分析它们之间的变化关系．
\par
解通过“五点法”画出函数3，一sin产3-工十
\par
SlnX
\par
4
\par
周期内的简图，如图5．4一7．
\par
y=sin@+却
\par
y=slnx
\par
y=sin(x-\ocrglyph{00E5})
\par
3烹
\par
2
\par
3丌
\par
5訌
\par
4
\par
2
\par
4
\par
4
\par
图5．4一7
\par
观察图5．4一7，可以发现：3，=sin，十了）的图象可以由3，一sin工的图象上每一
\par
点0，sinx）的纵坐标不变、横坐标减去匹得到，也就是将3，=sinx的图象向左平移
\par
玉个单位长度得到．
\par
类似地，观察图5，4一7，可以发现：3一、i"（'一一）的图象可以由y=sin工的图
\par
4
\par
象上每一点0，sin的纵坐标不变、横坐标加上匹得到，也就是将的图象
\par
向右平移匹个单位长度得到．
\par
一般地，函数3一sin(x十@（常数0）的图象可以由3一sin的图象向左（当
\par
>0）或向右（当<0）平移《到个单位长度得到．3，一sin(x十的值域为[一1，1]
\par
周期为2丌．
\end{ocrdraft}
\end{bankitem}
